% Part: turing-machines
% Chapter: machines-computations
% Section: variants

\documentclass[../../../include/open-logic-section]{subfiles}

\begin{document}

\olfileid{tur}{mac}{var}
\olsection{ટ્યુરિંગ મશીનોનાં વિવિધ સ્વરૂપો}

હકીકતમાં ટ્યુરિંગ મશીનોની વ્યાખ્યા આપવાની ઘણી રીતો છે;
આપણી રીત માત્ર એક છે. કેટલીક બાબતમાં આપણી વ્યાખ્યા
બીજી વ્યાખ્યાઓ કરતાં વધુ છૂટ આપે છે. આપણે મનસ્વી
સાન્ત વર્ણમાળાઓને મંજૂરી આપીએ છીએ; વધુ મર્યાદિત
વ્યાખ્યા ટેપનાં માત્ર બે પ્રતીકો, $\TMstroke$ અને~$\TMblank$,
મંજૂર કરી શકે. આપણે મશીનને ટેપ પર પ્રતીક લખતાં જ
ખસવાની છૂટ આપીએ છીએ; બીજી વ્યાખ્યાઓમાં લખવું
અથવા ખસવું, બેમાંથી એક જ થઈ શકે. આપણે ટેપના હેડને
ખસેડ્યા વગર લખવાની શક્યતા મંજૂર કરીએ છીએ; બીજી
વ્યાખ્યાઓ $\TMstay$ ``સૂચના'' છોડે છે. બીજી બાબતમાં
આપણી વ્યાખ્યા વધુ મર્યાદિત છે. આપણે ટેપ માત્ર એક
દિશામાં અનંત છે એમ ધાર્યું છે; બીજી વ્યાખ્યાઓ તેને ડાબે
અને જમણે બંને તરફ અનંત રાખે છે. હકીકતમાં જુદી જુદી
ગમે તેટલી ટેપ, અથવા ખાનાંની અનંત જાળી પણ મંજૂર કરી
શકાય. આપણે ટ્યુરિંગ મશીનનો સૂચનાઓનો ગણ સંક્રમણ
વિધેયથી દર્શાવીએ છીએ; બીજી વ્યાખ્યાઓ સંક્રમણ સંબંધ
વાપરે છે, જેમાં કોઈ પણ આપેલી સ્થિતિમાં મશીન પાસે
એકથી વધુ શક્ય સૂચનાઓ હોય છે.

વ્યાખ્યામાં કરેલી આ છેલ્લી છૂટ ખાસ રસપ્રદ છે. આપણી
વ્યાખ્યામાં મશીન અવસ્થા~$q$માં પ્રતીક~$\sigma$ વાંચે,
ત્યારે $\delta(q, \sigma)$ નવું પ્રતીક, નવી અવસ્થા અને
ટેપના હેડનું નવું સ્થાન નક્કી કરે છે. પરંતુ સૂચનાઓનો
ગણ વર્તમાન અવસ્થા-પ્રતીક જોડીઓ $\tuple{q,
  \sigma}$ અને નવી અવસ્થા-પ્રતીક-દિશા ત્રિકો
$\tuple{q', \sigma',
  D}$ વચ્ચેનો સંબંધ હોય, તો ટ્યુરિંગ મશીનની ક્રિયા
એક જ રીતે નક્કી થાય એવું જરૂરી નથી---સૂચનાઓના
સંબંધમાં $\tuple{q,
  \sigma, q', \sigma', D}$ અને
$\tuple{q, \sigma, q'', \sigma'', D'}$ બંને હોઈ શકે.
આ કિસ્સામાં આપણી પાસે \emph{અનિર્ધારિત} ટ્યુરિંગ મશીન
છે. સંગણકીય જટિલતાના સિદ્ધાંતમાં આવાં મશીનો અગત્યની
ભૂમિકા ભજવે છે.

ટ્યુરિંગ મશીન ક્યારે થંભે તેની પણ જુદી રીતો છે: આપણે
સંક્રમણ વિધેય અવ્યાખ્યાયિત હોય ત્યારે તે થંભે એમ કહીએ
છીએ; બીજી વ્યાખ્યાઓમાં મશીન ખાસ નિર્ધારિત થંભવાની
અવસ્થામાં હોય તે જરૂરી છે. આવી અવસ્થા જરૂરી રાખવાથી
ટ્યુરિંગ મશીનો શું ગણી શકે તેમાં ફેર પડતો નથી તે આપણે
\olref[hal]{sec}માં સમજાવ્યું છે. આપણા ટ્યુરિંગ મશીનોની
ટેપ માત્ર એક દિશામાં અનંત હોવાથી ક્યારેક મશીન સૂચના
યોગ્ય રીતે ચલાવી શકતું નથી: સૌથી ડાબું ખાનું વાંચતું
હોય અને ડાબે ખસવાનું કહેવામાં આવે ત્યારે. આપણી
વ્યાખ્યા મુજબ ``નીચે પડી જવા''ને બદલે તે એ જ સ્થાને
રહે છે; પરંતુ આવું બને ત્યારે મશીન થંભે એવી વ્યાખ્યા
પણ આપી શકીએ. આ વ્યાખ્યા પણ સમકક્ષ છે: ખાનું~$0$માંથી
ડાબે ખસવાનો પ્રયત્ન કરતાં થંભતું ટ્યુરિંગ મશીનનું વર્તન
અનુકરવા, દરેક સંક્રમણ $\delta(q,\TMendtape) =
\tuple{q',\sigma,\TMleft}$ દૂર કરી શકીએ---પછી
મશીન~$\TMendtape$ પર ડાબે ખસવાનો પ્રયત્ન કરવાને
બદલે થંભે છે.\footnote{આ રીત \emph{પૂરેપૂરી} કામ
કરતી નથી, કારણ કે ખાનું~$0$ સિવાયનાં ખાનાંમાં પણ
$\TMendtape$ લખતાં અને વાંચતાં આપણને કોઈ રોકતું નથી
(જુઓ \olref[una]{ex:mover}). આવાં કામ માટે બીજું
પ્રતીક~$\TMendtape'$ ઉમેરીને આ મુશ્કેલી દૂર કરી શકીએ.}

સંખ્યાઓ દર્શાવવાની પણ જુદી રીતો છે (અને તેથી ટ્યુરિંગ
મશીન ગણતું ઇનપુટ-આઉટપુટ વિધેય પણ બદલાય છે): આપણે
એક-પ્રતીકી નિરૂપણ વાપરીએ છીએ, પરંતુ દ્વિ-આધારી
નિરૂપણ પણ વાપરી શકાય. તે માટે $\TMblank$ અને~$\TMendtape$
ઉપરાંત બે પ્રતીકો જોઈએ.

હવે એક રસપ્રદ હકીકત: કયાં વિધેયો ટ્યુરિંગ-સંગણનીય છે
તેના પર આમાંથી કોઈ ફેરફારનો પ્રભાવ પડતો નથી.
\emph{એક વ્યાખ્યા અનુસાર કોઈ વિધેય ટ્યુરિંગ-સંગણનીય
હોય, તો બધી વ્યાખ્યાઓ અનુસાર તે ટ્યુરિંગ-સંગણનીય છે.}

આ વાતની ચકાસણીની વિગતોમાં નહીં જઈએ. માત્ર એક
ઉદાહરણ જોઈએ: ટેપ બંને દિશામાં અનંત હોય અથવા
અનેક ટેપ હોય એવી છૂટ આપવાથી સંગણનશક્તિ વધતી
નથી. લગભગ આવું કારણ છે: એક જ, એક દિશામાં
અનંત ટેપવાળું ટ્યુરિંગ મશીન અનેક ટેપનું અથવા બે
દિશામાં અનંત ટેપનું અનુકરણ કરી શકે છે. દાખલા તરીકે,
વધારાની અવસ્થાઓ અને સૂચનાઓ વાપરીને અનેક ટેપવાળા
અથવા બે દિશામાં અનંત ટેપવાળા મશીનના પ્રોગ્રામને
એક જ, એક દિશામાં અનંત ટેપવાળા મશીનના પ્રોગ્રામમાં
``રૂપાંતરિત'' કરી શકીએ. રૂપાંતરિત મશીન સમ ક્રમાંકનાં
ખાનાં ટેપ~$1$નાં ખાનાં (અથવા બે દિશામાં અનંત ટેપનાં
``ધન'' ખાનાં) માટે અને વિષમ ક્રમાંકનાં ખાનાં ટેપ~$2$
નાં ખાનાં (અથવા ``ઋણ'' ખાનાં) માટે વાપરી શકે છે.

\end{document}
